\section*{Chapitre \ref{ch:g10:functions} --- Fonctions}

\begin{solution}{exo:g10:functions:1}
$f(0) = 1$~; \quad $f(2) = 8 - 6 + 1 = 3$~; \quad
$f(-1) = 2 + 3 + 1 = 6$~; \quad
$f\!\left(\tfrac12\right) = 2 \times \tfrac14 - \tfrac32 + 1 = \tfrac12 -
\tfrac32 + 1 = 0$.
\end{solution}

\begin{solution}{exo:g10:functions:2}
$f$~: aucune contrainte, \omterm{def:g10:functions:function}{ensemble de définition} $\R$.

$g$~: exige $x + 4 \neq 0$, \omterm{def:g10:functions:function}{ensemble de définition} tous les réels sauf
$-4$.

$h$~: exige $2x - 6 \geq 0$, c'est-à-dire $x \geq 3$~: \omterm{def:g10:functions:function}{ensemble de
définition} $\intco{3}{+\infty}$.

$k$~: $x^2 + 1 \geq 1 > 0$ ne s'annule jamais, \omterm{def:g10:functions:function}{ensemble de définition}
$\R$.
\end{solution}

\begin{solution}{exo:g10:functions:3}
\emph{\omterm{def:g10:functions:function}{Antécédents} de $0$~:} $x^2 - 2x = x(x-2) = 0$, donc $0$ et $2$.

\emph{\omterm{def:g10:functions:function}{Antécédents} de $3$~:} $x^2 - 2x = 3$, c'est-à-dire
$x^2 - 2x - 3 = 0$, c'est-à-dire $(x-3)(x+1) = 0$ (vérifier en
développant), donc $3$ et $-1$.

\emph{\omterm{def:g10:functions:function}{Antécédents} de $-1$~:} $x^2 - 2x + 1 = (x - 1)^2 = 0$, donc
l'unique \omterm{def:g10:functions:function}{antécédent} $1$.
\end{solution}

\begin{solution}{exo:g10:functions:4}
$f(-1) = 1 + 2 + 1 = 4$~: oui, $A(-1, 4)$ est sur le \omterm{def:g10:functions:graph}{graphique}.
$f(2) = 4 - 4 + 1 = 1$~: oui, $B(2, 1)$ est aussi sur le \omterm{def:g10:functions:graph}{graphique}.
\end{solution}

\begin{solution}{exo:g10:functions:5}
\emph{1.} La plus grande valeur du tableau est $4 = g(3)$ (\omterm{def:g10:functions:extrema}{maximum}), la
plus petite est $-2 = g(0)$ (\omterm{def:g10:functions:extrema}{minimum}).

\emph{2.} Sur $\intcc{-3}{0}$ la \omterm{def:g10:functions:function}{fonction} décroît, et $-1 < -0.5$,
donc $g(-1) \geq g(-0.5)$.

\emph{3.} Sur chaque \omterm{def:g10:numbers:interval}{intervalle} de monotonie, $g$ prend chaque valeur au
plus une fois, donc $g(x) = 0$ a au plus une solution par \omterm{def:g10:numbers:interval}{intervalle}~: au
plus $3$ solutions au total. (Ici $0$ est entre $-2$ et $1$, entre $-2$ et
$4$, et entre $0$ et $4$, donc il y en a exactement trois.)
\end{solution}

\begin{solution}{exo:g10:functions:6}
Soit $u < v$. Alors
\[
f(v) - f(u) = (-2v + 7) - (-2u + 7) = -2(v - u) < 0,
\]
car $v - u > 0$. Donc $f(v) < f(u)$~: $f$ est strictement \omterm{def:g10:functions:variations}{décroissante}
sur $\R$.
\end{solution}

\begin{solution}{exo:g10:functions:7}
Compléter le carré~: $(x+3)^2 = x^2 + 6x + 9$, donc
\[
f(x) = x^2 + 6x + 11 = (x + 3)^2 + 2 .
\]
Pour tout $x$, $(x+3)^2 \geq 0$, donc $f(x) \geq 2$~; et
$f(-3) = 0 + 2 = 2$. Le \omterm{def:g10:functions:extrema}{minimum} de $f$ est $2$, atteint en $x = -3$.
\end{solution}

\begin{solution}{exo:g10:functions:8}
\emph{1.} Longueur $+$ largeur $= 20$, donc la longueur est $20 - x$ et
\[
A(x) = x(20 - x).
\]
Les deux dimensions doivent être positives~: $0 < x < 20$.

\emph{2.} $A(4) = 64$, $A(8) = 96$, $A(10) = 100$, $A(12) = 96$,
$A(16) = 64$. Les valeurs croissent jusqu'à $x = 10$ puis décroissent
symétriquement~: l'aire semble maximale pour $x = 10$, c'est-à-dire pour
un jardin \emph{carré}. (\cref{ch:g11:quad} prouve cette conjecture.)
\end{solution}

\begin{solution}{exo:g10:functions:9}
\emph{1.} Pour tout $x$~: $x^2 + 1 \geq 1 > 0$, donc
$f(x) = \frac{1}{x^2+1}$ est positif, et comme $x^2 + 1 \geq 1$, en
prenant les inverses (les deux membres positifs) on obtient
$f(x) \leq 1$.

\emph{2.} \omterm{def:g10:functions:extrema}{Maximum}~: $f(0) = 1$ et $f(x) \leq 1$ pour tout $x$, donc $1$
est le \omterm{def:g10:functions:extrema}{maximum}, atteint en $0$. \omterm{def:g10:functions:extrema}{Minimum}~: $f(x) > 0$ toujours, mais
aucune valeur de $x$ n'atteint $0$ (un quotient de nombres non nuls est
non nul), et $f$ prend des valeurs aussi proches de $0$ que l'on veut
pour $x$ grand~; donc $f$ n'a pas de \omterm{def:g10:functions:extrema}{minimum}.
\end{solution}

\begin{solution}{exo:g10:functions:10}
\emph{1.} $f(v) - f(u) = v^2 - u^2 = (v-u)(v+u)$. Si $0 \leq u < v$,
alors $v - u > 0$ et $v + u > 0$ (somme d'un nombre positif ou nul et
d'un nombre positif), donc $f(v) - f(u) > 0$~: $f$ est strictement
\omterm{def:g10:functions:variations}{croissante} sur $\intco{0}{+\infty}$.

\emph{2.} Si $u < v \leq 0$, alors $v - u > 0$ encore, mais maintenant
$v + u < 0$ (somme d'un nombre négatif ou nul et d'un nombre négatif),
donc $f(v) - f(u) < 0$~: $f$ est strictement \omterm{def:g10:functions:variations}{décroissante} sur
$\intoc{-\infty}{0}$.
\end{solution}

\begin{solution}{pb:g10:functions:1}
\textbf{1.} \omterm{def:g10:functions:function}{Ensemble de définition}~: $x \neq 0$ (division par $x$).
$f(1) = \frac12(1 + 2) = \frac32$~; $f(2) = \frac12(2 + 1) =
\frac32$~; $f\!\left(\frac32\right) = \frac12\left(\frac32 +
\frac43\right) = \frac{17}{12}$.

\textbf{2.} $\frac12\left(x + \frac2x\right) = \frac32$ donne
$x^2 + 2 = 3x$, c'est-à-dire $x^2 - 3x + 2 = (x - 1)(x - 2) = 0$~:
les \omterm{def:g10:functions:function}{antécédents} de $\frac32$ sont $1$ et $2$ --- comme la question 1
l'annonçait.

\textbf{3.} $f(x) = x$ donne $x + \frac2x = 2x$, donc $\frac2x = x$
et $x^2 = 2$~: \omterm{pb:g10:functions:1}{points fixes} $\sqrt2$ et $-\sqrt2$. Le \omterm{pb:g10:functions:1}{point fixe}
positif est la diagonale du carré unité, irrationnelle d'après le
devoir du week-end sur l'irrationalité du volume précédent.

\textbf{4.} $1 \to \frac32 \to \frac{17}{12} \to \frac{577}{408}$~:
les \omterm{def:g10:numbers:approx}{approximations} de $\sqrt2$ dues à Héron. Sa recette «~faire la
moyenne de l'estimation et de $2/$estimation~» est exactement
\emph{l'itération de la \omterm{def:g10:functions:function}{fonction} $f$} --- et le nombre que le procédé
poursuit est le \omterm{pb:g10:functions:1}{point fixe} de $f$.

\textbf{5.} Les \omterm{pb:g10:functions:1}{points fixes} sont les \omterm{def:g10:numbers:interunion}{intersections} de la courbe avec
la droite $y = x$~: réinjecter les sorties fait cheminer le point le
long de la courbe vers cette \omterm{def:g10:numbers:interunion}{intersection}. Dans le devoir du week-end
sur les deux thermomètres, du volume précédent, la lecture
$-40^\circ$ reposait sur la même idée de croisement avec la
diagonale, pour la \omterm{def:g10:functions:function}{fonction} de conversion.

\textbf{6.} $f(v) - f(u) = \frac{v - u}{2} + \frac1v - \frac1u
= \frac{v - u}{2} + \frac{u - v}{uv}
= (v - u)\left(\frac12 - \frac{1}{uv}\right)
= \frac{(v - u)(uv - 2)}{2uv}$. Pour $0 < u < v \leq \sqrt2$~:
$uv < 2$, la parenthèse est négative, donc $f(v) < f(u)$~: $f$ est
\omterm{def:g10:functions:variations}{décroissante}. Pour $\sqrt2 \leq u < v$~: $uv > 2$, donc $f$ est
\omterm{def:g10:functions:variations}{croissante}. Point de retournement~: $\sqrt2$.

\textbf{7.} \omterm{def:g10:functions:variations}{Décroissante} avant $\sqrt2$, \omterm{def:g10:functions:variations}{croissante} après~: $f$
atteint son \omterm{def:g10:functions:extrema}{minimum} en $\sqrt2$, et ce \omterm{def:g10:functions:extrema}{minimum} vaut
$f(\sqrt2) = \frac12\left(\sqrt2 + \frac{2}{\sqrt2}\right) =
\sqrt2$. Conséquence~: toute sortie de la machine (pour une entrée
positive) vaut au moins $\sqrt2$ --- dès le premier tour de
manivelle, les estimations vivent dans $\intco{\sqrt2}{+\infty}$.

\textbf{8.} Si $x > \sqrt2$, alors $x^2 > 2$, donc $\frac2x < x$ et
la moyenne $f(x)$ de $x$ et de $\frac2x$ est inférieure à $x$~; et
$f(x) \geq \sqrt2$ d'après la question 7, avec égalité seulement en
$x = \sqrt2$. Ainsi $\sqrt2 < f(x) < x$~: chaque itération améliore
strictement sans jamais dépasser --- la suite de la question 4
descend en glissant sur $\sqrt2$.

\textbf{9.} Sur $\intoo{0}{+\infty}$~: $f$ décroît depuis $+\infty$
(au voisinage de $0$) jusqu'au \omterm{def:g10:functions:extrema}{minimum} $\sqrt2$ atteint en
$x = \sqrt2$, puis croît sans borne.

\textbf{10.} $g(x) = x$ donne $x^2 = x$, soit $x(x - 1) = 0$~: \omterm{pb:g10:functions:1}{points
fixes} $0$ et $1$. À partir de $0.9$~: $0.81$, $0.656$, $0.430$,
$0.185$ --- glissement vers $0$. À partir de $1.1$~: $1.21$,
$1.464$, $2.144$, $4.595$ --- fuite vers l'infini. Le \omterm{pb:g10:functions:1}{point fixe} $0$
attire, le \omterm{pb:g10:functions:1}{point fixe} $1$ repousse~: un cheveu d'écart au départ,
des destins opposés.

\textbf{11.} $7 \to 22 \to 11 \to 34 \to 17 \to 52 \to 26 \to
13 \to 40 \to 20 \to 10 \to 5 \to 16 \to 8 \to 4 \to 2 \to 1$~:
seize étapes, avec un sommet à $52$.

\textbf{12.} $27 \to 82 \to 41 \to 124 \to 62 \to 31 \to 94 \to
47 \to 142 \to 71 \to 214 \to \dots$ --- toujours en train de monter
après dix étapes, en route vers un sommet à $9\,232$ et un vol de
$111$ étapes. Leçon~: une règle de deux lignes peut produire un
comportement que personne ne devine au premier coup d'œil --- des
nombres voisins ($26$, $27$, $28$) ont des vols radicalement
différents.

\textbf{13.} Une puissance de $2$ redescend toute son échelle par
divisions successives~: $2^k \to 2^{k-1} \to \dots \to 2 \to 1$. En
bas, $h(1) = 4$, $h(4) = 2$, $h(2) = 1$~: le vol entre dans le cycle
$4 \to 2 \to 1 \to 4$. Un \omterm{pb:g10:functions:1}{point fixe} exigerait $\frac n2 = n$
(seulement $n = 0$) ou $3n + 1 = n$ (négatif)~: il n'y en a aucun
parmi les entiers strictement positifs --- les vols observés
s'achèvent non pas en un \omterm{pb:g10:functions:1}{point fixe}, mais dans ce petit \emph{cycle}.

\textbf{14.} Un contre-exemple serait soit un vol qui croît
indéfiniment (sans jamais revenir à $1$), soit un second cycle
disjoint de $4 \to 2 \to 1$. Une vérification jusqu'à $10^{20}$
laisse une infinité de nombres non testés --- exactement la leçon des
régions du cercle, dans le devoir du week-end sur l'oracle aux
allumettes du volume précédent~: un accord sur un nombre fini de cas
ne démontre rien sur tous.

\textbf{15.} Pour la machine de Héron, nous avions une
\emph{structure} --- une différence factorisée (question 6) révélant
les variations, un \omterm{def:g10:functions:extrema}{minimum}, un encadrement --- tandis que la règle à
bascule de parité des grêlons n'offre aucune prise de ce genre~:
personne n'a trouvé la structure qui la dompterait.

\textbf{16.} $\sqrt{x - 3}$ exige $x - 3 \geq 0$~:
$\intco{3}{+\infty}$. $\frac{1}{x^2 - 9}$ exige $x \neq \pm 3$~:
$\R$ privé de $-3$ et de $3$. $\sqrt{9 - x^2}$ exige $x^2 \leq 9$~:
$\intcc{-3}{3}$.

\textbf{17.} $f(x) = 3$ donne $x + \frac2x = 6$, donc
$x^2 - 6x + 2 = 0$, c'est-à-dire $(x - 3)^2 = 7$~:
$x = 3 - \sqrt7$ ou $x = 3 + \sqrt7$ (tous deux positifs~: deux
\omterm{def:g10:functions:function}{antécédents} exacts).

\textbf{18.} La formule a un sens physique tant que la pierre est en
l'air~: $H(t) = 5t(4 - t) \geq 0$ pour $t \in \intcc{0}{4}$. Forme
canonique~: $H(t) = 20 - 5(t - 2)^2$~: \omterm{def:g10:functions:extrema}{maximum} $20$ m à $t = 2$ s.
Variations~: \omterm{def:g10:functions:variations}{croissante} sur $\intcc{0}{2}$ de $0$ à $20$,
\omterm{def:g10:functions:variations}{décroissante} sur $\intcc{2}{4}$ jusqu'à $0$.

\textbf{19.} $20t - 5t^2 \geq 15 \iff t^2 - 4t + 3 \leq 0 \iff
(t - 1)(t - 3) \leq 0 \iff t \in \intcc{1}{3}$~: deux secondes
pleines au-dessus de $15$ m.

\textbf{20.} Chaque machine possède un \omterm{def:g10:functions:function}{ensemble de définition} (là où
la règle est autorisée), une règle (la formule ou la bascule de
parité) et une courbe ou un tableau qui l'affiche tout entière. À
chacune nous avons posé les deux grandes questions~: \emph{où mène la
répétition} (\omterm{pb:g10:functions:1}{points fixes}, cycles, attraction --- réglé pour Héron,
ouvert pour les grêlons) et \emph{comment la sortie bouge avec
l'entrée} (variations, extremums --- le sommet de la pierre, le
\omterm{def:g10:functions:extrema}{minimum} de Héron). Les chapitres à venir affinent les deux~: les
\omterm{def:g10:functions:function}{fonctions} de référence d'abord, puis, les années suivantes, les
dérivées pour mesurer les variations et les limites pour certifier
où les itérations aboutissent.
\end{solution}
